% Source: 01 - Mon  Sep 28 2026.pdf, page 10. Content remains in source order.
\sourcepage{01}{10}

\section{Abstract Decision Problems}
\[
\Pi_D\subseteq\mathcal I\times\mathcal S,\qquad\mathcal S=\{\mathrm{yes},\mathrm{no}\}.
\]
Fundamental parts in the formulation of an abstract decision problem:
\begin{itemize}
\item description of the generic instance.
\item question on the instance.
\item the solution is a yes/no answer to the question.
\end{itemize}

\subsection{EXAMPLE: Membership of integer key in a set}

\[
\left\{\begin{aligned}
\text{INSTANCE: }&\langle S,m\rangle:\ S\subseteq\mathbb N\text{, finite},\ m\in\mathbb N,\\
\text{QUESTION: }&m\in S?
\end{aligned}\right.
\]
 
A decision problem is a function
\[
\Pi_D:\mathcal I\longrightarrow\{\mathrm{yes},\mathrm{no}\}.
\]
Yes/No instances are also called Positive/negative instances.

\subsection{Restriction to decision problems: excessive?}
We restrict our analysis on computational problems to decision problems. This is not a severe restriction since many optimization problems can be reduced to decision problems.
In practice we usually solve more complex optimization problems.
